\section{Xử lý song song trong .NET/C\#}

\subsection{Cơ chế data parallelism}

\textbf{Data parallelism} là áp dụng cùng một phép tính lên nhiều phần tử dữ liệu.
\textbf{Task Parallel Library (TPL)} là thư viện của .NET hỗ trợ xử lý song song.
\textbf{ThreadPool} là tập thread được .NET quản lý và tái sử dụng, giúp lập trình
viên không phải tự tạo và thu hồi từng thread. Khi một vòng lặp song song bắt đầu,
TPL chia nguồn thành các
\textbf{partition} (những nhóm phần tử nhỏ hơn), đưa các nhóm lên ThreadPool và chờ
toàn bộ nhóm hoàn tất \cite{ms-tpl,ms-threadpool,ms-data-parallelism}.

Trong \figref{fig:parallel-flow}, mỗi ``worker'' là một ThreadPool thread đang thực hiện một partition.
Số worker thực tế do môi trường thực thi .NET và cấu hình giới hạn mức độ song song
quyết định.

\begin{center}
\resizebox{0.80\textwidth}{!}{\begin{tikzpicture}[
  node distance=7mm and 10mm,
  every node/.style={font=\small,align=center},
  flowbox/.style={draw,rounded corners=2mm,minimum height=10mm,minimum width=23mm,thick},
  arrow/.style={-{Latex[length=2mm]},thick}
]
\node[flowbox,fill=gray!12,draw=gray!65] (source) {Dữ liệu\\đầu vào};
\node[flowbox,fill=blue!12,draw=blue!55,right=of source] (split) {Chia thành\\các partition};
\node[flowbox,fill=blue!15,draw=blue!55,above right=of split] (w1) {ThreadPool\\worker 1};
\node[flowbox,fill=green!15,draw=green!55!black,right=of split] (w2) {ThreadPool\\worker 2};
\node[flowbox,fill=orange!18,draw=orange!70!black,below right=of split] (w3) {ThreadPool\\worker 3};
\node[flowbox,fill=gray!12,draw=gray!65,right=15mm of w2] (result) {Chờ và ghép\\kết quả};
\draw[arrow] (source) -- (split);
\draw[arrow] (split) -- (w1);
\draw[arrow] (split) -- (w2);
\draw[arrow] (split) -- (w3);
\draw[arrow] (w1) -- (result);
\draw[arrow] (w2) -- (result);
\draw[arrow] (w3) -- (result);
\end{tikzpicture}
}
\captionof{figure}{Luồng chia, thực hiện và ghép kết quả của một vòng lặp Parallel}
\label{fig:parallel-flow}
\end{center}

Việc chia và ghép công việc không miễn phí. \textbf{Overhead} là thời gian phụ dùng
để chia dữ liệu, lập lịch, đồng bộ và ghép kết quả. Nếu mỗi phần tử chỉ cần một phép
tính rất nhanh, overhead có thể lớn hơn thời gian tiết kiệm được. Vì vậy Parallel
chỉ có cơ hội nhanh hơn khi lượng tính toán đủ lớn và các phần ít phụ thuộc nhau
\cite{ms-plinq-speedup}.

\subsection{Các API xử lý nhiều phần tử}

\textbf{Collection} là tập dữ liệu trong bộ nhớ như mảng hoặc danh sách;
\textbf{iteration} là một lần thực hiện thân vòng lặp. Mỗi API dưới đây phù hợp với
một dạng dữ liệu và cách xử lý khác nhau.
\textbf{LINQ} là nhóm cú pháp và API dùng để truy vấn dữ liệu trong C\#;
\textbf{PLINQ} là phiên bản thực thi song song một số thao tác LINQ trên collection;
\texttt{Parallel.ForEachAsync} dành cho hàm bất đồng bộ và có từ .NET 6, nên bản
chất khác các vòng lặp Parallel đồng bộ \cite{ms-parallel-api}.

{\small
\begin{longtable}{|>{\raggedright\arraybackslash}p{0.22\textwidth}|
                  >{\raggedright\arraybackslash}p{0.30\textwidth}|
                  >{\raggedright\arraybackslash}p{0.38\textwidth}|}
\caption{Các API xử lý nhiều phần tử thường dùng}\label{tab:parallel-api}\\
\hline
\textbf{API} & \textbf{Phù hợp khi} & \textbf{Điểm cần chú ý} \\
\hline
\texttt{Parallel.For} & Dữ liệu được truy cập bằng chỉ số và mỗi phần tử cần tính toán CPU độc lập. & Là API đồng bộ; chỉ trả về khi toàn bộ vòng lặp hoàn tất. \\
\hline
\texttt{Parallel.ForEach} & Cần áp dụng cùng một phép tính CPU-bound lên từng phần tử của collection. & Không dùng để bọc các lời gọi I/O đồng bộ. \\
\hline
PLINQ với \texttt{AsParallel()} & Chuỗi LINQ có bước lọc, biến đổi hoặc tổng hợp đủ nặng. & Có chi phí chia, ghép dữ liệu và giữ thứ tự kết quả. \\
\hline
\texttt{Parallel.\allowbreak ForEachAsync} & Nhiều thao tác bất đồng bộ cần giới hạn số thao tác đang hoạt động. & Không đồng nghĩa với xử lý CPU song song; API này có từ .NET 6. \\
\hline
\end{longtable}
}

\subsection{Thay đổi từ .NET Core 3.1 đến .NET 8}

{\small
\begin{longtable}{|>{\raggedright\arraybackslash}p{0.29\textwidth}|
                  >{\raggedright\arraybackslash}p{0.29\textwidth}|
                  >{\raggedright\arraybackslash}p{0.32\textwidth}|}
\caption{Thay đổi trực tiếp của Parallel từ .NET Core 3.1 đến .NET 8}\label{tab:parallel-version}\\
\hline
\textbf{.NET Core 3.1} & \textbf{.NET 8} & \textbf{Ý nghĩa khi sử dụng} \\
\hline
Đã có \texttt{Parallel.For}, \texttt{Parallel.ForEach} và PLINQ. & Các API cốt lõi và cách tư duy data parallelism vẫn được giữ. & Code cũ thường không cần đổi chỉ vì chuyển sang phiên bản .NET mới. \\
\hline
Chưa có \texttt{Parallel.ForEachAsync}. & .NET 8 có API này vì nó được bổ sung từ .NET 6. & Có sẵn API để chạy hàm bất đồng bộ trên collection, nhận tín hiệu dừng và giới hạn số thao tác đồng thời. \\
\hline
\end{longtable}
}

Thay đổi trực tiếp dễ thấy nhất là \texttt{Parallel.ForEachAsync}. Việc nâng phiên bản .NET
không bảo đảm mọi vòng lặp nhanh hơn: tốc độ còn phụ thuộc dữ liệu, phần cứng, tải hệ
thống và cách viết code, nên vẫn phải đo trên trường hợp thực.
